\documentclass[11pt]{article}

\usepackage[margin=1in]{geometry}
\usepackage{amsmath,amssymb}
\usepackage{booktabs}
\usepackage{hyperref}
\usepackage{microtype}
\usepackage{xcolor}

\hypersetup{
  colorlinks=true,
  linkcolor=blue,
  urlcolor=blue
}

\newcommand{\VOtwo}{VO\textsubscript{2}}

\title{Practical Comparison of the Quasistatic and Dynamic-Phase \VOtwo{} Neuristor Models}
\author{Gabriel Berezovsky}
\date{\today}

\begin{document}
\maketitle

\section{Purpose}

This note explains the practical difference between:
\begin{enumerate}
    \item the quasistatic \VOtwo{} resistance model used in the original Yuanhang Zhang repository; and
    \item a local experimental extension that gives the \VOtwo{} phase fraction a finite response time.
\end{enumerate}

The distinction is important because the sustained current-source oscillations found in the recent parameter search depend on the second model. They are not a direct reproduction of the original Yuanhang Zhang mathematics.

\section{Shared Physical Ingredients}

Both models describe a thermally active \VOtwo{} resistor. The electrical resistance depends strongly on temperature because \VOtwo{} changes between insulating and metallic phases.

The resistance law is
\begin{equation}
R(T,g) = R_m + R_0 \exp\left(\frac{E_a/k_B}{T}\right) g,
\label{eq:resistance}
\end{equation}
where:
\begin{itemize}
    \item $R_m$ is the metallic resistance floor;
    \item $R_0\exp[(E_a/k_B)/T]$ is the insulating contribution; and
    \item $g$ is the insulating phase fraction, with approximately $g=1$ for insulating material and $g=0$ for metallic material.
\end{itemize}

The hysteresis model calculates an equilibrium phase fraction
\begin{equation}
g_{\mathrm{eq}} = g_{\mathrm{eq}}(T,\mathcal{H}),
\label{eq:geq}
\end{equation}
where $\mathcal{H}$ represents the heating/cooling branch and reversal-history variables.

For the ideal current-source circuit used in this project, the electrical and thermal equations are
\begin{align}
C\frac{dV}{dt} &= I_{\mathrm{in}}-\frac{V}{R}, \label{eq:current}\\
C_{\mathrm{th}}\frac{dT}{dt} &= \frac{V^2}{R}-S_e(T-T_0). \label{eq:thermal}
\end{align}

The current source supplies current to the parallel combination of the capacitor and \VOtwo{} device. Joule heating raises the device temperature, while $S_e(T-T_0)$ removes heat.

\section{Original Quasistatic Model}

\subsection{Mathematics}

In the original Yuanhang Zhang model, the phase fraction is calculated directly from the instantaneous temperature and hysteresis state:
\begin{equation}
g(t) = g_{\mathrm{eq}}\!\left(T(t),\mathcal{H}(t)\right).
\label{eq:quasistatic}
\end{equation}

Therefore, the resistance used at each timestep is
\begin{equation}
R(t) = R\!\left(T(t),g_{\mathrm{eq}}(T(t),\mathcal{H}(t))\right).
\end{equation}

There is no separate differential equation for $g$. The independent dynamical states are only voltage and temperature:
\begin{equation}
\mathbf{x}_{\mathrm{quasi}} =
\begin{bmatrix}
V \\ T
\end{bmatrix}.
\end{equation}

\subsection{Practical Meaning}

``Quasistatic'' means that the material phase is assumed to respond immediately compared with the electrical and thermal timescales. If temperature changes, the phase fraction and resistance immediately move to the value prescribed by the hysteresis curve.

For example, if heating changes the equilibrium insulating fraction from $0.8$ to $0.2$, the model immediately uses $g=0.2$. It does not model the finite time needed for metallic domains to nucleate or grow.

\subsection{Consequence Under an Ideal Current Source}

At a steady state,
\begin{align}
V_* &= I_{\mathrm{in}}R(T_*),\\
I_{\mathrm{in}}^2R(T_*) &= S_e(T_*-T_0).
\end{align}

Our searches found that the quasistatic ideal-current system can produce switching transients and damped ringing, but it eventually settles to such a fixed point for the tested fitted resistance law and physical parameter ranges.

\section{Local Dynamic-Phase Extension}

\subsection{Mathematics}

The local extension treats the actual phase fraction $g(t)$ as an additional dynamical state. It approaches the hysteresis-prescribed equilibrium fraction with a finite time constant $\tau_g$:
\begin{equation}
\tau_g\frac{dg}{dt}
=
g_{\mathrm{eq}}(T,\mathcal{H})-g.
\label{eq:dynamic_g}
\end{equation}

The resistance is then calculated from the delayed actual phase fraction:
\begin{equation}
R(t) = R\!\left(T(t),g(t)\right).
\end{equation}

The dynamical state vector becomes
\begin{equation}
\mathbf{x}_{\mathrm{dynamic}} =
\begin{bmatrix}
V \\ T \\ g
\end{bmatrix}.
\end{equation}

The quasistatic model is recovered in the limiting case
\begin{equation}
\tau_g \rightarrow 0,
\qquad
g(t)\rightarrow g_{\mathrm{eq}}(T(t),\mathcal{H}(t)).
\end{equation}

\subsection{Practical Meaning}

The dynamic model says that the material does not transform instantaneously. Temperature first changes what phase fraction is energetically preferred, but the actual material state takes time to follow.

For example, suppose heating changes the equilibrium fraction from $0.8$ to $0.2$. With finite $\tau_g$, the actual $g$ may still be near $0.8$ for several nanoseconds and then gradually approach $0.2$.

This delay changes the resistance, voltage, Joule heating, and subsequent temperature trajectory.

\section{Why the Delay Can Produce Oscillation}

The dynamic-phase feedback loop is:
\begin{enumerate}
    \item Current causes Joule heating.
    \item Temperature rises and changes the preferred equilibrium phase fraction $g_{\mathrm{eq}}$.
    \item The actual phase fraction $g$ responds late.
    \item Because $g$ is delayed, resistance and voltage remain temporarily inconsistent with the current temperature.
    \item The delayed resistance change modifies Joule heating and can cause cooling.
    \item During cooling, the phase fraction again responds late.
    \item The delayed heating/cooling feedback can repeat as a limit cycle.
\end{enumerate}

The key point is that $\tau_g$ introduces a new memory and delay. A delayed negative-feedback system can oscillate even when the corresponding instantaneous-feedback system settles.

\section{Direct Comparison}

\begin{center}
\begin{tabular}{p{0.27\textwidth}p{0.31\textwidth}p{0.31\textwidth}}
\toprule
Property & Quasistatic model & Dynamic-phase extension \\
\midrule
Independent states & $V,T$ & $V,T,g$ \\
Phase fraction & $g=g_{\mathrm{eq}}(T,\mathcal{H})$ & $\tau_g\dot g=g_{\mathrm{eq}}-g$ \\
Phase response & Immediate & Delayed \\
New parameter & None & $\tau_g$ \\
Original Yuanhang code & Yes & No \\
Current-source result in our search & Transient ringing, then fixed point & Sustained oscillatory domains found \\
\bottomrule
\end{tabular}
\end{center}

\section{What Can Be Claimed}

\subsection{Valid statement for the original model}

For the fitted specimen resistance law and tested physical parameter ranges, the original quasistatic ideal-current-source model did not produce a robust sustained oscillator. It produced transient switching or ringing before settling.

\subsection{Valid statement for the dynamic extension}

Adding a finite phase-response time can create sustained current-source oscillations. This demonstrates that delayed phase-transition kinetics are a possible mechanism for oscillation.

\subsection{Statement that should not be made yet}

The dynamic-extension results should not be presented as a reproduction of Yuanhang Zhang's original model. The extension must be independently justified and calibrated before its quantitative predictions are treated as experimental predictions.

\section{What Would Validate the Dynamic Extension}

The parameter $\tau_g$ should ideally be supported by one or more of:
\begin{itemize}
    \item time-resolved measurements of the \VOtwo{} phase transition;
    \item measured delay between temperature, voltage, and resistance changes;
    \item literature values for domain nucleation and growth times;
    \item comparison against raw oscilloscope traces over several currents; or
    \item a more detailed physical phase-field or kinetic model.
\end{itemize}

Until then, $\tau_g$ is best described as a physically motivated phenomenological parameter.

\section{Summary}

The quasistatic model assumes:
\begin{equation}
g = g_{\mathrm{eq}}(T,\mathcal{H}).
\end{equation}

The dynamic extension assumes:
\begin{equation}
\tau_g\dot g = g_{\mathrm{eq}}(T,\mathcal{H})-g.
\end{equation}

This apparently small change adds a third dynamical state and a delayed feedback loop. That delay is what allowed sustained oscillations to appear in the ideal current-source simulations. The mechanism is physically plausible, but it is a local experimental extension rather than part of the original Yuanhang Zhang repository.

\end{document}
